problem:  
Let \( (M^{4n+3}, g, \mathcal{S} = \{(\xi_i,\eta_i,\Phi_i)\}_{i=1}^{3}) \) be a **compact quasi‑regular 3‑Sasakian manifold**, with locally free \(\mathrm{Sp}(1)\)-action generated by the Reeb vector fields and orbifold quaternionic‑Kähler base \( B = M / \mathrm{Sp}(1)\).  
Let  
\[
\rho : \mathrm{Aut}(M,\mathcal{S}) \longrightarrow \mathrm{Iso}(B)
\]
be the natural projection induced by the quotient map.

**Open problem (precise conjecture):**  
Is the kernel of \(\rho\) always isomorphic to \(\mathrm{Sp}(1)\) for every compact quasi‑regular 3‑Sasakian manifold?  

Equivalently:  
Do orbifold singularities in \(B\) ever enlarge the Reeb‑generated subgroup inside \(\mathrm{Aut}(M,\mathcal{S})\)?  
In other words, can one find a quasi‑regular 3‑Sasakian manifold for which  
\[
\ker(\rho) \supsetneq \mathrm{Sp}(1)
\]
as Lie groups, implying the existence of automorphisms acting trivially on the base but nontrivially on the total space, not generated by the standard Reeb action?

Why is it a "good" problem:  
This formulation isolates a **specific, unresolved structural question** in 3‑Sasakian geometry. In the regular (smooth) case, the Reeb fields give an exact \(\mathrm{Sp}(1)\)-kernel of \(\rho\), and all known examples conform to this pattern. However, in the quasi‑regular setting, orbifold singularities in the quaternionic‑Kähler base can introduce hidden isotropy phenomena. Whether these can **create extra vertical automorphisms** remains completely unstudied.  

The problem, therefore, probes the fine interplay between **orbifold topology** and **symmetry structure** in special Riemannian geometries—a delicate bridge connecting Sasakian geometry, Lie group actions, and the theory of quaternionic‑Kähler orbifolds. Its statement is simple, yet its resolution would fundamentally clarify how base singularities influence total‑space symmetries, potentially leading to new constructions or rigidity results in Sasakian and Einstein geometries. It exemplifies the “narrow entrance, profound depth” characteristic of a truly good mathematical problem.